% Part: second-order-logic
% Chapter: metatheory
% Section: undecidability-and-axiomatizability

\documentclass[../../../include/open-logic-section]{subfiles}

\begin{document}

\olfileid{sol}{met}{nax}

\olsection{દ્વિતીય-ક્રમ તર્કશાસ્ત્ર સ્વયંસિદ્ધીકરણીય નથી}


\begin{thm}
\ollabel{thm:sol-undecidable}
દ્વિતીય-ક્રમ તર્કશાસ્ત્ર અનિર્ણેય છે.
\end{thm}

\begin{proof}
પ્રથમ-ક્રમ !!{sentence} પ્રથમ-ક્રમ તર્કશાસ્ત્રમાં પ્રમાણભૂત ત્યારે
અને માત્ર ત્યારે જ હોય જ્યારે તે દ્વિતીય-ક્રમ તર્કશાસ્ત્રમાં
પ્રમાણભૂત હોય; વળી પ્રથમ-ક્રમ તર્કશાસ્ત્ર અનિર્ણેય છે.
\end{proof}

\begin{thm}
\ollabel{cor:sol-not-axiomatizable}
દ્વિતીય-ક્રમ તર્કશાસ્ત્ર માટે અસરકારક રીતે આપેલું એવું કોઈ યથાર્થ અને
પૂર્ણ !!{derivation} તંત્ર નથી.
\end{thm}
\sourcecorrection{OLSOL-008}{સ્થિર અંગ્રેજી મૂળમાં અહીં કોઈપણ યથાર્થ
અને પૂર્ણ નિષ્પત્તિ-તંત્ર નથી એમ નિઃશરત કહેવાયું છે. અસરકારક રીતે
આપેલું તંત્ર હોવાની શરત જરૂરી છે: અસરકારકતાની શરત વિના બધાં પ્રમાણભૂત
વાક્યોને જ સ્વયંસિદ્ધિઓ તરીકે લઈ શકાય. નીચેની સાબિતી સંગણનીય
પરિગણનાનો વિરોધાભાસ ઉત્પન્ન કરે છે, તેથી આ શરત સ્પષ્ટ કરી છે.}

\begin{proof}
$!A$ અંકગણિતની ભાષાનું !!a{sentence} હોય. $\Sat{N}{!A}$ ત્યારે
અને માત્ર ત્યારે જ જ્યારે $\Th{PA^2} \Entails !A$. $!P$ એ
$\Th{PA^2}$નાં નવ સ્વયંસિદ્ધિઓનો સંયોજક હોય. $\Th{PA^2} \Entails !A$
ત્યારે અને માત્ર ત્યારે જ જ્યારે $\Entails !P \lif
!A$, એટલે દરેક સંરચના માટે
$\Sat{M}{!P \lif !A}$. 
\sourcecorrection{OLSOL-009}{સ્થિર અંગ્રેજી મૂળના આ સંતોષ-સૂત્રમાં
બીજી \texttt{\textbackslash Sat} દલીલનું બંધ કૌંસ ગણિત-મર્યાદાની
બહાર મૂકાયું છે, તેથી મૂળ TeX અસંતુલિત છે. અહીં
$\Sat{M}{!P \lif !A}$નું બંધ કૌંસ સૂત્રની અંદર મૂક્યું અને
પહેલાંના પ્રામાણ્ય-દાવાને અનુરૂપ દરેક સંરચનાનો વ્યાપ સ્પષ્ટ કર્યો છે.}
હવે~$\Obj{0}$ની જગ્યાએ~$z$, $\prime$ની જગ્યાએ એકસ્થાની
વિધેય-ચલ~$u$, $+$ અને~$\times$ની જગ્યાએ અનુક્રમે દ્વિસ્થાની
વિધેય-ચલો~$u'$ અને~$u''$, તથા $<$ની જગ્યાએ દ્વિસ્થાની સંબંધ-ચલ~$L$
મૂકીને, પછી આ બધા ચલો પર સર્વવાચી પરિમાણીકરણ કરતાં મળતું શુદ્ધ
દ્વિતીય-ક્રમ !!{sentence}
$\lforall[z][\lforall[u][\lforall[u'][\lforall[u''][\lforall[L][(!P' \lif !A')]]]]]$
વિચારો. આ વાક્ય પ્રમાણભૂત ત્યારે અને માત્ર ત્યારે જ જ્યારે મૂળ વાક્ય
પ્રમાણભૂત હોય; તે ત્યારે અને માત્ર ત્યારે જ જ્યારે
$\Th{PA^2} \Entails !A$, અને તે ત્યારે અને માત્ર ત્યારે જ જ્યારે
$\Sat{N}{!A}$. તેથી દ્વિતીય-ક્રમ તર્કશાસ્ત્ર માટે યથાર્થ અને પૂર્ણ
અસરકારક સાબિતી-તંત્ર હોત, તો તેની સાબિતીઓની પરિગણનામાંથી આ ખાસ
આકારનાં વાક્યો પસંદ કરીને~$\Struct{N}$માં સાચાં !!{sentence}sની
સંગણનીય પરિગણના~$f\colon \Nat \to \Sent[L_A]$ મેળવી શકાત. આ
વિધેય~$\Th{Q}$માં કોઈ પ્રથમ-ક્રમ સૂત્ર~$!B_f(x, y)$ વડે નિરૂપણક્ષમ
હોત. ત્યારે !!{formula}~$\lexists[x][!B_f(x, y)]$ એ
$\Struct{N}$માં સાચાં પ્રથમ-ક્રમ !!{sentence}sની ગોડેલ-સંખ્યાઓનો
ગણ વ્યાખ્યાયિત કરત, જે ટાર્સ્કીના પ્રમેયનો વિરોધ કરે છે.
\sourcecorrection{OLSOL-010}{સ્થિર અંગ્રેજી મૂળ વાક્યોની સંગણનીય
પરિગણનાથી અંકગણિતમાં વ્યાખ્યેય ગણ સુધીના પગલામાં વાક્યોની સંખ્યાકીય
સંકેતીકરણ-પ્રક્રિયા સ્પષ્ટ કરતું નથી. અહીં $y$ને સાચા વાક્યની
ગોડેલ-સંખ્યા તરીકે ઓળખાવી છે; વાક્ય પોતે અંકગણિતના ચલનું મૂલ્ય નથી.}
\end{proof}



\end{document}
